\documentclass[10pt]{article}
\usepackage[T1]{fontenc}
\usepackage[utf8]{inputenc}
\usepackage[french]{babel}
\usepackage{amsmath,amssymb,mathrsfs,mathpazo}
\usepackage{geometry}
\geometry{a4paper,margin=2cm}
\usepackage{graphicx,array,tabularx,tabularray,colortbl,multirow,multicol,booktabs,enumitem,pifont,ulem,fancybox,diagbox,cancel}
\usepackage{pstricks,pstricks-add,pst-plot,pst-tree,pst-node,pst-text,pst-func,pst-math,pst-bspline,pst-3dplot,pst-eucl,pst-solides3d}
\let\pscalculate\undefined
\usepackage{pst-bezier}
\usepackage{tikz,pgfplots,tkz-tab}
\usetikzlibrary{arrows,arrows.meta,patterns,calc,positioning,trees,decorations.markings,shapes.geometric,shapes.misc,shapes.arrows,intersections,decorations.pathmorphing,decorations.pathreplacing,backgrounds,shadows,babel}
\pgfplotsset{compat=1.18}
\usepgfplotslibrary{fillbetween}
\usepackage[np]{numprint}
\usepackage{listings,makecell,pas-tableur,twoopt}
\let\barre\undefined
\usepackage[pstricks]{tablvar}
\definecolor{darkgreen}{RGB}{0,100,0}
\providecommand{\red}{\color{red}}
\providecommand{\blue}{\color{blue}}
\providecommand{\green}{\color{green}}
\providecommand{\black}{\color{black}}
\providecommand{\white}{\color{white}}
\providecommand{\gray}{\color{gray}}
\providecommand{\cyan}{\color{cyan}}
\providecommand{\magenta}{\color{magenta}}
\providecommand{\yellow}{\color{yellow}}
\providecommand{\violet}{\color{violet}}
\providecommand{\orange}{\color{orange}}
\providecommand{\euro}{\texteuro}
\providecommand{\EUR}{\texteuro}
\providecommand{\ieme}{\textsuperscript{e}}
\providecommand{\degres}{\ensuremath{^\circ}}
\providecommand{\degree}{\ensuremath{^\circ}}
\providecommand{\vv}[1]{\overrightarrow{#1}}
\providecommand{\Coord}[1]{#1}
\providecommand{\decosix}{}
\providecommand{\starredbullet}{\ensuremath{\bullet}}
\pagestyle{empty}
\setlength{\parindent}{0pt}
\usepackage{adjustbox}\DeclareUnicodeCharacter{2236}{:}\DeclareUnicodeCharacter{2190}{\ensuremath{\leftarrow}}
\providecommand{\R}{\mathbb{R}}
\providecommand{\N}{\mathbb{N}}
\providecommand{\Z}{\mathbb{Z}}
\providecommand{\Q}{\mathbb{Q}}
\providecommand{\D}{\mathbb{D}}
\providecommand{\C}{\mathbb{C}}
\providecommand{\e}{{\mathrm{e}}}
\providecommand{\vect}[1]{\overrightarrow{#1}}
\providecommand{\vectt}[1]{\overrightarrow{\mathrm{#1}}}
\providecommand{\vvt}[1]{\overrightarrow{\mathrm{#1}}}
\providecommand{\np}[2][]{{#2~#1}}
\providecommand{\vv}[1]{{\overrightarrow{#1}}}
\renewcommand{\d}{{\mathrm{d}}}
\providecommand{\degres}{{^\circ}}
\providecommand{\degre}{{^\circ}}
\providecommand{\blue}{\color{blue}}
\providecommand{\red}{\color{red}}
\providecommand{\enskip}{\quad}
\providecommand{\llbracket}{\lbrack\!\lbrack}
\providecommand{\rrbracket}{\rbrack\!\rbrack}
\providecommand{\euro}{{\text{€}}}
\providecommand{\no}{{\text{n°}}}
\providecommand{\interval}[2]{{[#1;#2]}}
\providecommand{\intervFO}[2]{{[#1;#2[}}
\providecommand{\dotfill}{\ldots}

\providecommand{\ds}{\displaystyle}
\providecommand{\textding}[1]{#1}
\providecommand{\Oij}{\(\mathrm{O};\vec{\imath},\vec{\jmath}\)}
\providecommand{\Oijk}{\(\mathrm{O};\vec{\imath},\vec{\jmath},\vec{k}\)}
\providecommand{\Ouv}{\(\mathrm{O};\vec{u},\vec{v}\)}
\providecommand{\intervFF}[2]{\ensuremath{[#1;#2]}}
\renewcommand{\euro}{\ifmmode\text{\texteuro}\else\texteuro\fi}
\DeclareUnicodeCharacter{00B0}{\ifmmode^{\circ}\else\textdegree\fi}
\DeclareUnicodeCharacter{20AC}{\ifmmode\text{\texteuro}\else\texteuro\fi}
% Énoncés : Éducation nationale. Transcription et corrigés : APMEP, Denis Vergès et contributeurs. Découpage et classement : CoopMaths, Sébastien Lozano, Éric Elter et contributeurs. Adaptation de la présentation : Maths Explorer.
% CC BY-SA : https://coopmaths.fr/www/about/
% https://coopmaths.fr/alea/static/bac/2022/tex/bac_2022_05_sujet1_ameriquenord_1_cor.tex
\begin{document}
\noindent\textit{APMEP / CoopMaths. CC BY-SA. Présentation adaptée par Maths Explorer.}\par\medskip
\medskip
\textbf{Partie A}

\medskip
\begin{enumerate}
	\item 
	\begin{enumerate}
		\item L'arbre complété avec les valeurs disponibles :
		\begin{center}
				\pstree[treemode=R,nodesepA=0pt,nodesepB=2.5pt,treesep=1cm,levelsep=2.5cm]{\TR{}}
{\pstree{\TR{$V$~}\taput{$\frac23$}}
{\TR{$R$}\taput{$\frac{1}{50}$}
\TR{$\overline{R}$}\tbput{$\frac{49}{50}$}
}
\pstree{\TR{$\overline{V}$~}\tbput{$\frac13$}}
{\TR{$R$}\taput{$\frac{1}{10}$}
	\TR{$\overline{R}$}\tbput{$\frac{9}{10}$}
}
}
\end{center}
		\item D'après la formule des probabilités totales :
		
		$p(R)=p(R\cap V)+ p\left(R \cap \overline{V}\right) = p_V(R)\times p(V) + p_{\overline{V}}(R)\times p\left(\overline{V}\right) = \dfrac{1}{50}\times \dfrac23+\dfrac{1}{10}\times \dfrac13 = \dfrac{7}{150}$
		\item On cherche la probabilité : $p_R(V)$.
		
D'après la formule de Bayes, $p_R(V)=\dfrac{p(R\cap V)}{p(R)}=\dfrac{\dfrac{1}{50}\times \dfrac23}{\dfrac{7}{150}}=\dfrac{2}{7}$.
	\end{enumerate}
	\item 
	\begin{enumerate}
		\item Il s'agit là d'un schéma de Bernoulli : la répétition de 20 expériences aléatoires n'ayant que deux issues, identiques et indépendantes entre elles. $X$ suit donc une loi binomiale de paramètres $n = 20$ et $p = \dfrac23$ : $X \sim \mathcal{B}\left(20~;~\dfrac23\right)$
		\item $p(X=10)=\dbinom{20}{10}\times \left(\dfrac23\right)^5 \times \left(1-\dfrac23\right)^{20-10}\approx 0,054$.
		\item $p(X\geqslant 10)=1-p(X<10)=1-p(X\leqslant 9)\approx 0,962$
		\item Calculons l'espérance $E(X)$ de cette loi binomiale: $E(X)=n\times p=20 \times \dfrac23=\dfrac{40}{3}\approx 13,333$

En moyenne, il se rendra à la gare en vélo 13,33 jours par mois soit 13 jours à l'unité près.
	\end{enumerate}
	\item Calculons l'espérance de la loi de probabilité de $T$ :
	
	$\begin{gathered}E(t)=0,14\times 10+0,13 \times 11+0,13\times12+0,12\times13+0,12\times14+0,11\times15+0,10\times16+0,08\times17+0,07\times18\\
	=13,5\end{gathered}$
	
	Il mettra en moyenne 13 minutes et demie  pour se rendre à la gare en voiture.
\end{enumerate}

\pagebreak

\end{document}
